% Implementation, Oversight, Enforcement, and Fiscal Provisions. Four movements;
% the dense numbers are carried in the implementation-choices table and the
% fiscal table. All figures are planning indices or simulation results.

This section sets out how the bill is implemented, who oversees and enforces it,
how it is phased in, and how it is paid for. The implementation levers are
mapped in Table~\ref{tab:impl-levers} and the fiscal case in
Table~\ref{tab:impl-fiscal}; the prose synthesises the design rather than
narrating the tables.

\subsection{Implementation Design Choices}

From the standard menu of implementation levers, the bill selects a coherent set
that all serve the verification-before-generation rule. It requires mandatory
disclosure of artificial-intelligence use in any robot-patient interaction; it
secures auditability through the hash-chained audit trail and the five-server
MCP-PAI topology, with its role-based access and alignment to the
electronic-records rule \cite{kawchak2026mcp, cfr-part11}; it retains human review
of the medical-necessity decision; it provides a Food and Drug Administration
review touchpoint consistent with the device and credibility framework
\cite{fda2023credibility}; it opens a reimbursement pathway for verified Physical
AI procedures under the Medicare standard \cite{usc-1395y}; it seeds site
readiness through a grant program; and it imposes a periodic reporting duty. The
infrastructure and governance choices draw on the national platform's MCP and
implementation-strategy sections, including its standards-integration gates and
its quality and risk management. Table~\ref{tab:impl-levers} records each lever
and its source.

\begin{table}[ht]
\centering
\begin{tabularx}{\textwidth}{L{3.4cm} Y}
\toprule
\textbf{Lever} & \textbf{Bill choice and source} \\
\midrule
\multicolumn{2}{l}{\textit{Parent: physical-ai .../final\_paper/sections/}} \\
Disclosure & Mandatory disclosure of AI use in robot-patient interactions \\
Auditability & Hash-chained trail, five MCP servers, 23 tools (national\_mcp.tex) \\
Human review & Retained medical-necessity decision (regulatory\_landscape.tex) \\
FDA touchpoint & Device and credibility review (fda2023credibility) \\
Reimbursement & CMS pathway for verified procedures (financial\_analysis.tex) \\
Grant program & Seeds site readiness toward the PSL gate (implementation\_strategy.tex) \\
Reporting duty & Periodic performance and USL reporting (implementation\_strategy.tex) \\
\bottomrule
\end{tabularx}
\caption{Implementation levers selected by the bill and their sources.}
\label{tab:impl-levers}
\end{table}

\subsection{Oversight and Enforcement Authority}

Oversight is split along the lines the existing law already draws. A Food and
Drug Administration touchpoint reviews the device and credibility dimensions,
while a state insurance or health authority reaches the
utilization-review-adjacent duties. The enforcement mechanics are audits and
compliance reviews, civil penalties for willful violations on the California
model \cite{ca-sb1120}, a commissioner finding of noncompliance that triggers
production of the underlying algorithms on the Texas model \cite{tx-sb1822}, and
certification of bias minimization on the New York model \cite{ny-a9149}. The
operative enforcement trigger throughout is a single fact: the absence of a
cleared verification-before-generation record for a robot-patient interaction.
Where a denial or a halt is involved, the Florida human-review backstop applies,
so that no automated decision stands as the sole basis for action without
qualified human review \cite{fl-hb527}.

\subsection{Phased Rollout and Workforce}

The bill follows a three-phase national rollout. The first phase activates a
single state, building on the California legislative authorization and the
operational record of the 24-hour, 168-patient simulation that demonstrated 99.7
percent uptime with zero harm. The second phase expands to five to ten sites
across multiple states under the five-server MCP architecture, with site
selection governed by the Physical AI Standard Level and robot admission by the
Unification Standard Level minimum. The third phase reaches nationwide coverage
across many sites, with governance transitioning from centralized to federated.
The workforce roles, Physical AI system operators, robot maintenance technicians,
and safety specialists, and their certification under the adapted good-clinical-
practice training standard, are drawn from the national platform's
implementation-strategy timeline and workforce tables and are referenced rather
than reproduced \cite{kawchak2026national, kawchak2026sitedocs}. Site activation
is tied to the Physical AI Standard Level gate and robot admission to the
Unification Standard Level minimum, so the rollout and the readiness gates of
Section~\ref{sec:statutory_text} advance together.

\subsection{Budget, Funding, and Offsets}

The fiscal case is summarised in Table~\ref{tab:impl-fiscal}; every figure is a
planning index or a simulation result, not a guarantee. The central offset
argument is that automating the VVUQ assurance layer is inexpensive relative to
its decision weight: an autonomous agent carried the heavier assurance half of
the work cheaply and reproducibly, as the cost evidence of the two studies shows,
so the gate that makes the safety decision is among the least expensive
components of the trial \cite{kawchak2026vvuq01, kawchak2026vvuq02}. Against that
modest assurance cost stands the cost of delay, where each day of slower drug
development is materially expensive, so the savings from faster and safer trials
are the funding offset for the grant program and the oversight costs. The 2026
real-time clinical trials initiative reinforces the case, because automated
assurance is what makes a real-time trial feasible at all \cite{fda2026realtime}.

\begin{table}[ht]
\centering
\begin{tabularx}{\textwidth}{L{4.4cm} Y}
\toprule
\textbf{Fiscal item} & \textbf{Planning index or simulation result} \\
\midrule
\multicolumn{2}{l}{\textit{Source: financial\_analysis.tex; figures are planning ranges}} \\
Cost of trial delay & Each day of delayed approval is materially costly (driver of the offset) \\
Per-patient cost & Lower than traditional trials through continuous operation and automated capture \\
Phase 1 capital (1 to 2 sites) & About \$15.5M to \$24.5M (California activation) \\
Phase 2 capital (5 to 10 sites) & About \$72M to \$120M (multi-state) \\
Phase 3 capital (50 to 100 sites) & About \$377M to \$630M (nationwide) \\
Assurance-layer cost & Inexpensive relative to its decision weight (the offset argument) \\
Patient and sponsor benefit & Faster access; more trials within existing budgets \\
\bottomrule
\end{tabularx}
\caption{Fiscal case as planning indices and simulation results, not guarantees.}
\label{tab:impl-fiscal}
\end{table}
